\begin{tabularx}{\textwidth}{L{24mm}X}
\toprule
\textbf{Szempont} & \textbf{Értékelés} \\
\midrule
Előny & Gyors visszacsatolást ad, jobban kezeli a bizonytalan követelményeket, és a rendszer a felhasználói tapasztalatok alapján finomítható. \\
Előny & A specifikáció inkrementálisan fejlődhet, ezért az ügyfél hamarabb láthat működő részeket, mint egy szigorúan szekvenciális modellben. \\
Hátrány & A folyamat kevésbé átlátható a menedzsment számára, mert nem mindig vannak jól elkülöníthető, lezárt dokumentum-alapú mérföldkövek. \\
Hátrány & A sok módosítás ronthatja a rendszer szerkezetét, ha nincs tudatos refaktorálás, architekturális kontroll és dokumentálás. \\
Tipikus alkalmazás & Rövidebb élettartamú, kis és közepes rendszereknél, felhasználói felületeknél, webes rendszereknél, illetve olyan projektekben, ahol a követelmények csak használat közben tisztulnak. \\
\bottomrule
\end{tabularx}
